\documentclass[11pt]{article}
\usepackage[T1]{fontenc}
\usepackage[utf8]{inputenc}
\usepackage{lmodern}
\usepackage[margin=0.84in]{geometry}
\usepackage{amsmath,amssymb,amsthm,mathtools}
\usepackage{booktabs,longtable,array}
\usepackage[hidelinks]{hyperref}
\usepackage{xcolor}
\usepackage{enumitem}
\usepackage{microtype}
\usepackage{url}

\newtheorem{definition}{Definition}
\newtheorem{theorem}{Theorem}
\newtheorem{proposition}{Proposition}
\newtheorem{corollary}{Corollary}
\newcommand{\R}{\mathbb R}
\newcommand{\RobInv}{\mathsf{RobInv}_{\mathrm{src}}}
\newcommand{\Vad}{\mathcal V_{\mathrm{adm}}}
\newcommand{\EqSrc}{\mathsf{EqSrc}_{\mathrm{prot}}}
\newcommand{\status}[1]{\textsf{#1}}
\newcolumntype{P}[1]{>{\raggedright\arraybackslash}p{#1}}
\setlist{nosep,leftmargin=*}
\urlstyle{tt}

\title{Source-Law-Space Robust-Invariance Protection\\
\large V22 P4--T02 B2 proposal-only formalization}
\author{Ontology Formalizer parent synthesis, RT-20260811-007}
\date{11 August 2026}

\begin{document}
\maketitle

\begin{center}
\fbox{\begin{minipage}{0.94\linewidth}
\textbf{Decisive result:}
\status{theorem\_with\_hypotheses\_and\_proof}.
For one explicitly typed regular finite-barrier source-law datum, uniform
strict conormal margins for every admitted generator and variation are
equivalent to strict tangent inclusion and imply strong forward invariance on
the declared solution domain.  The formal relation and its noncanonical
certificate orbit descend through coherent \(C^1\) source-presentation
arrows.  A balanced pair of admitted normal variations at a zero-margin face
gives an exact obstruction.

\smallskip
\textbf{Status and stop.}
The roots \(X_{\mathrm{law}},K,\mathcal A,\Delta\), and \(\EqSrc\) are
disclosed \emph{proposal-only source-extension data}.  This theorem does not
derive or adopt them from current ontology.  Exactly one focused Smuggling
Auditor packet is selected but not executed.  All five local freezes remain;
all fourteen Distance-to-GR rows remain \status{no\_delta}; D7 is not
reevaluated; B2 is inactive; and P4--T03 is locked.
\end{minipage}}
\end{center}

\section{Authorized question and fixed roots}

Handoff-1020 authorizes exactly one formalization-or-obstruction packet:
\begin{flushleft}
\small\ttfamily
PKT-V22-P4T02-B2-SOURCE-LAW-SPACE-ROBUST-INVARIANCE-\\
PROTECTION-FORMALIZATION-V1.
\end{flushleft}
The missing premise isolated by RT005 and selected by RT006 is not another
token, quotient, common character, oriented-matroid classifier, or special
cubic flow.  It is a law on a full source-state arena that specifies which
states, generators, and variations are admissible and protects that set under
every admitted variation.

This task therefore predeclares the tuple
\[
  \mathcal D_{\mathrm{prot}}
   =(X_{\mathrm{law}},K,\mathcal A,\Delta,\EqSrc)
\]
as new proposal-only source-extension data.  The tuple is fixed before the
protection calculation.  Its predeclaration prevents adaptive alteration in
this packet; it does not prove that current ontology selects or derives the
tuple.  That epistemic question belongs to the separately admitted audit.

No target atlas, metric, tetrad, coframe, desired GR cone, physical norm,
detector response, empirical result, or benchmark calibration is an input.
The source-law chart is not adopted \(M_{\mathrm{src}}\), its tangent cones
are not physical characteristic cones, and its presentation arrows are not
asserted to be physical gauge.

\section{Typed source-law datum}

\begin{definition}[Regular proposal-only protection datum]
Let \(X_{\mathrm{law}}\) be a finite-dimensional real \(C^1\) chart.  Let
\(K\subset X_{\mathrm{law}}\) be closed and full-dimensional.  Near every
relevant boundary point it has a finite presentation
\[
 K=\{x:\rho_i(x)\geq0,\ i=1,\ldots,m\},                 \tag{1}
\]
where the active differentials are linearly independent.  Let \(\Theta\) be
compact and let \(\mathcal A=\{F_\theta\}_{\theta\in\Theta}\) be a
predeclared, jointly continuous, locally equi-Lipschitz family of source
generators.  Let \(\Delta(x)\) be a predeclared, nonempty compact-valued,
locally bounded, outer-semicontinuous variation multifunction, with the same
properties inherited by \(\Vad\).  Its admitted selections satisfy the stated
Carathéodory hypotheses.  Finally, let \(\EqSrc\) be a groupoid of justified
locally \(C^{1,1}\) diffeomorphisms, coherent parameter homeomorphisms, and
locally \(C^{1,1}\) inverse arrows for the complete tuple.
\end{definition}

For \(x\in\partial K\), write
\(I(x)=\{i:\rho_i(x)=0\}\).  The regular tangent and positive-conormal
objects are
\begin{align}
 T_K(x)&=\{v:d\rho_i(x)[v]\geq0\ \text{for every }i\in I(x)\}, \tag{2}\\
 N_K^+(x)&=\operatorname{pos}\{d\rho_i(x):i\in I(x)\}.          \tag{3}
\end{align}
The second object uses inward conormals; it is not the signed physical normal
to a causal cone.

Define the admitted velocity multifunction
\[
 \Vad(x)=\bigcup_{F\in\mathcal A}\bigl(F(x)+\Delta(x)\bigr).
\]

\begin{definition}[Strict source-law-space robust invariance]
\begin{equation}
 \RobInv(K,\mathcal A,\Delta)
 \quad\Longleftrightarrow\quad
 \Vad(x)\subset\operatorname{Int}T_K(x)
 \quad\text{for every relevant }x\in\partial K .
 \label{eq:robinv}
\end{equation}
The result is \status{fail} when an exact violating velocity is supplied and
\status{undefined} when a root, an effective boundary test, or its coherent
transport is absent.  Missing data never count as a pass.
\end{definition}

\section{Tangent characterization and forward invariance}

\begin{theorem}[Regular active-conormal characterization]
\label{thm:tangent}
Under the regularity in (1), the contingent tangent cone is (2), and
\[
 \operatorname{Int}T_K(x)
 =\{v:d\rho_i(x)[v]>0\ \text{for every }i\in I(x)\}.
\]
\end{theorem}

\begin{proof}
The linear-independence constraint qualification identifies the contingent
cone of the regular inequality set with its linearized feasible cone.  This
is (2).  A finite intersection of independent closed half-spaces is
full-dimensional; its ambient interior is obtained by making every active
inequality strict.
\end{proof}

\begin{theorem}[Strict source-law-space strong invariance]
\label{thm:invariance}
Assume the regular datum above.  Assume every declared solution is absolutely
continuous, remains in a compact local solution tube on each bounded time
interval, and satisfies
\(\dot x(t)\in\Vad(x(t))\) almost everywhere.  If \(\Vad\) is locally bounded
and outer-semicontinuous and \eqref{eq:robinv} holds with locally uniform
positive active-conormal margins,
then every solution beginning in \(K\) remains in \(K\) throughout its
declared forward domain.
\end{theorem}

\begin{proof}
By Theorem~\ref{thm:tangent}, every active barrier derivative is positive for
every admitted boundary velocity.  Local boundedness and outer
semicontinuity of the velocity multifunction, continuity of the active
differentials, and the locally uniform margin preserve a positive half-margin
on a neighborhood of each contacted face.  On a bounded time interval the
contacted boundary portion is compact, so finitely many such neighborhoods
cover it.  If a solution had a first exit, at least one active
\(\rho_i(x(t))\) would have to cross from nonnegative to negative while its
absolutely continuous derivative is positive almost everywhere throughout
the last face neighborhood.  Integration gives an increase, contradicting
the crossing.  The argument applies simultaneously to every active face and
every admitted generator and variation.
\end{proof}

The strict margin is stronger than the weakest possible viability condition,
but it is deliberate: it supplies the perturbative protection that was absent
from the frozen signed-cubic candidate.  The theorem is conditional on the
declared solution and regularity class and makes no claim outside it.

\section{Quantitative certificates and their orbit}

\begin{theorem}[Finite conormal certificate]
\label{thm:certificate}
Suppose the barriers in (1) only present an independently supplied \(K\).
If, for every relevant boundary point, every active \(i\), every
\(F\in\mathcal A\), and every \(E\in\Delta(x)\), there are predeclared
numbers \(\alpha_i(x)>\delta_i(x)\geq0\) with
\begin{align}
 d\rho_i(x)[F(x)]&\geq\alpha_i(x), \\
 |d\rho_i(x)[E]|&\leq\delta_i(x),
\end{align}
then every active residual is at least
\(\alpha_i-\delta_i>0\), so
Theorem~\ref{thm:invariance} applies.
\end{theorem}

\begin{proof}
For every admitted sum,
\[
 d\rho_i[F+E]
 \geq d\rho_i[F]-|d\rho_i[E]|
 \geq\alpha_i-\delta_i>0.
\]
Theorems~\ref{thm:tangent} and \ref{thm:invariance} finish the proof.
\end{proof}

The barriers are certificates, not canonical source laws.  If, near an active
face,
\[
 \widetilde\rho_i=h_i\rho_{\pi(i)},\qquad h_i>0,
\]
then on that face
\[
 d\widetilde\rho_i=h_i\,d\rho_{\pi(i)}.
\]
Thus positive regraduation and face permutation scale the nominal bound,
variation bound, and residual margin together.  Redundant presentations are
retained as transported certificate fibres.  The output is an orbit or
groupoid-valued certificate, never a canonically selected raw barrier.

Positive regraduation preserves an already chosen inward orientation.  It
does not select that orientation or justify \(K\).

\section{Source-presentation transport}

\begin{theorem}[\(C^{1,1}\) presentation transport]
\label{thm:transport}
Let \((g,\eta_g)\) be a declared presentation arrow, where
\(g:X_{\mathrm{law}}\to X'_{\mathrm{law}}\) is locally \(C^{1,1}\) and
\(\eta_g:\Theta\to\Theta'\) is a homeomorphism.  Transport the complete
datum by
\begin{align*}
 K'&=gK,\\
 F'_{\eta_g(\theta)}(gx)&=Dg_xF_\theta(x),\\
 \Delta'(gx)&=Dg_x\Delta(x).
\end{align*}
Then
\[
 T_{K'}(gx)=Dg_xT_K(x),\qquad
 \Vad'(gx)=Dg_x\Vad(x),
\]
and \(\RobInv\) has the same pass, fail, or undefined value in both
presentations.  Declared solutions correspond by \(g\).
\end{theorem}

\begin{proof}
The contingent-cone chain rule gives the first identity.  The second is the
definition of complete pushforward.  Because \(Dg_x\) is an invertible linear
open map, it maps the interior of \(T_K(x)\) to the interior of
\(T_{K'}(gx)\).  Hence strict inclusion is equivalent.  The chain rule maps
absolutely continuous solutions forward, and the inverse arrow maps them
back.  Overlap coherence makes the result independent of a composable arrow
factorization.
\end{proof}

This is mathematical presentation naturality.  No physical gauge statement
follows.

\section{Exact obstruction at zero margin}

\begin{proposition}[Balanced-normal obstruction]
\label{prop:balanced}
Let \(x\in\partial K\), let \(\ell\in N_K^+(x)\) be a nonzero active inward
conormal, and suppose \(\ell(F(x))=0\).  If \(\Delta(x)\) contains \(E\) and
\(-E\) with \(\ell(E)\neq0\), then strict robust invariance fails at \(x\).
One admitted velocity is outside even the weak tangent cone.
\end{proposition}

\begin{proof}
The two admitted conormal velocities are
\[
 \ell(F+E)=\ell(E),\qquad
 \ell(F-E)=-\ell(E).
\]
They are opposite and nonzero.  One is negative, so its velocity violates
the corresponding weak tangent inequality in Theorem~\ref{thm:tangent}.
\end{proof}

The proposition is a falsifier for supplied data, not a global no-go theorem.
If the source-provenanced variation class contains the balanced pair, the
strict law fails.  Removing the pair only after seeing this failure would
trigger the countermodel-exclusion branch; removal requires independent
source provenance.

\section{Exact rational controls}

The positive fixture is the proposal-only orthant
\[
 K=\R_{\geq0}^3,
 \qquad
 F_\theta(x)=b_\theta-\operatorname{diag}(1,2,3)x,
\]
with
\[
 b_0=(3,4,5),\quad b_1=(4,5,6),\quad
 \Delta=[-1,1]\times[-2,2]\times[-3,3].
\]
On every active face the worst nominal residuals for \(F_0\) are exactly
\((2,2,2)\); for \(F_1\) they are \((3,3,3)\).  Every corner is covered
because the active inequalities are componentwise.

Regraduating the barriers by \((2,3,5)\) gives residuals
\((4,6,10)\).  The positive monomial arrow
\[
 g(x_1,x_2,x_3)=(2x_2,3x_1,5x_3)
\]
transports the same family to exact residuals \((4,6,10)\) and preserves
strictness.  The worst-case scalar
equilibria for the first generator are \((2,1,2/3)\), all positive.

For the balanced control, take \(x=(0,1,1)\),
\(\ell(v)=v_1\), a tangent nominal velocity, and the admitted pair
\(\pm(1,0,0)\).  The exact signs are \(+1\) and \(-1\).  A zero-margin
fixture produces no strict certificate, and omitting the provenance of
\(\Delta\) returns \status{undefined}, not pass.

The task-local exact arithmetic model reports twelve of twelve checks passed.
These fixtures validate equations and fail branches only.  They do not count
as source evidence, ontology adoption, physical evidence, or benchmark
recovery.

\section{Selection, meaning, and authority boundaries}

The map \(\RobInv\) classifies a supplied tuple.  It does not select the
tuple.  In particular:
\begin{itemize}
  \item a pass value does not determine \(K\), \(\mathcal A\), \(\Delta\),
  an orientation, or a presentation arrow;
  \item predeclaration and hashing establish non-adaptivity within this task,
  not an epistemic derivation from current source ontology;
  \item a source-law tangent cone is not a physical characteristic cone;
  \item forward invariance is not occurrence, probability, detector
  response, universal matter propagation, or empirical adequacy; and
  \item certificate naturality is not physical gauge or adoption authority.
\end{itemize}

The construction factors through none of the five frozen candidates.  It has
no shared \(\tau\), response-ray quotient, common character, holonomy datum,
oriented-matroid selector, binary sign token, half-line, or cubic generator.
The following freezes remain exact:
\begin{itemize}
  \item \path{NDCL-V22-P4T02-B2-SHARED-TAU-SELECTOR};
  \item \path{NDCL-V22-P4T02-B2-REPAIRED-QUOTIENT-DESCENT-ROBUSTNESS};
  \item \path{NDCL-V22-P4T02-B2-COMMON-CHARACTER-HOLONOMY-PROTECTION};
  \item \path{NDCL-V22-P4T02-B2-ORIENTED-MATROID-BRIDGE-SELECTION-ROBUSTNESS};
  \item \path{NDCL-V22-P4T02-B2-SIGNED-CUBIC-VIABILITY-SELECTOR-ROBUSTNESS}.
\end{itemize}

\section{Obligation and failure disposition}

All fourteen packet obligations are discharged at the proposal-only theorem
scope.  The source provenance obligations for \(K\), \(\mathcal A\), and
\(\Delta\) pass only as disclosed fixed roots pending audit; no
current-ontology derivation is claimed.  Regularity and solution hypotheses
are explicit; every declared stratum is covered; certificate and presentation
transport are proved; the balanced obstruction is proved; the result is
exactly one allowed theorem; and all downstream locks are retained.

The eleven fail branches remain executable guards.  Goal-selected \(K\),
post-outcome filtering of \(\mathcal A\), countermodel-selected narrowing of
\(\Delta\), or target-norm input would fail the candidate.  Exact tangent
violation returns fail; balanced zero-margin data trigger
Proposition~\ref{prop:balanced}; raw-certificate selection is replaced by an
orbit; incoherent transport fails; frozen-route factorization is absent;
physical overread is blocked; and missing or ineffective inputs return
undefined.

\section{Distance-to-GR matrix}

\begin{longtable}{P{0.27\linewidth}P{0.63\linewidth}}
\toprule
Burden & Status after RT007 \\
\midrule
\endfirsthead
\toprule
Burden & Status after RT007 \\
\midrule
\endhead
Source ontology primitives & \(X_{\mathrm{law}},K,\mathcal A,\Delta\), and
\(\EqSrc\) are proposal-only roots, not adopted primitives.  No delta.\\
Source equivalence \(\mathsf{EqSrc}\) & Conditional \(C^1\) transport is
proved for the proposed tuple; current source equivalence and physical gauge
are not derived.  No delta.\\
\(\mathsf{RetainH}\) & No retention or history primitive is supplied.  No
delta.\\
\(\mathsf{GenH}\) & A proposal-only generator family is an input, not a
current-ontology derivation.  No delta.\\
\(\mathsf{ObsLoc}_{lc}\) & No observer-localization or operational readout
map is supplied.  No delta.\\
\(\mathsf{Resp}_{lc}\) & Formal invariance is not detector or empirical
response semantics.  No delta.\\
\(M_{\mathrm{src}}\) & The chart is not adopted physical source spacetime.
No delta.\\
\(g_{\mathrm{eff}}\) & A source-law tangent cone is not conformal geometry,
physical scale, or an effective metric.  No delta.\\
Matter coupling & No sector-universal propagation or coupling law is derived.
No delta.\\
Einstein equations & No action, stress energy, variation principle, or field
equation is supplied.  No delta.\\
Finite-variation robustness & A conditional proposal-only theorem exists,
but source purity and current-ontology selection are unaudited.  No delta.\\
Benchmark promotion & No exact-GR comparison, recovery, or promotion occurs.
No delta.\\
Gate Chair status & No protected review, adoption, or verdict occurs.  No
delta.\\
Route freeze or hard-fail & All five local freezes remain active and no global
freeze is added.  No delta.\\
\bottomrule
\end{longtable}

\section{Parent-child synthesis and next packet}

The Physicist--Mathematician child supplies the regular domains, tangent and
positive-conormal formulas, strong-invariance proof, certificate orbit,
\(C^1\) transport theorem, balanced-normal obstruction, exact rational
controls, and theorem-only scope.  The Physicist--Philosopher child supplies
the distinction between predeclaration and epistemic derivation, orientation
nonselection, separation of ontology from mathematical model, and separation
of physical, empirical, adoption, and workflow status.  That child also
supplies smuggling falsifiers and the protected-authority boundary.  Both
select the same decisive
result and report no blocking conflict.  The parent retains the stricter LICQ
regular domain as an explicit theorem hypothesis; this narrows the theorem and
does not import target structure.

The positive result selects exactly one future packet:
\begin{flushleft}
\small\ttfamily
PKT-V22-P4T02-B2-SOURCE-LAW-SPACE-ROBUST-INVARIANCE-\\
PROTECTION-SMUGGLING-AUDIT-V1.
\end{flushleft}
Its type is \path{source_extension_smuggling_audit}; its role is
\path{smuggling-auditor@0.2.0}; its status is
\status{selected\_not\_executed}.  It must audit only whether the disclosed
roots and their orientation, margins, and arrows are genuinely target-free
and independently source-provenanced.  It may not treat this theorem or its
validation as adoption evidence.

\section{Explicit non-conclusions}

This task does not:
\begin{itemize}
  \item derive or adopt the candidate roots or law from current ontology;
  \item execute the selected Smuggling Auditor or any later Refuter packet;
  \item reopen or repair any frozen candidate;
  \item derive occurrence, probability, detector response, physical
  causality, time orientation, sector universality, conformal geometry,
  physical scale, \(g_{\mathrm{eff}}\), coupling, or field equations;
  \item reevaluate D7 or P2--T01 adequacy, act on B2, complete P4--T02 for
  plan dependency, or unlock P4--T03;
  \item prove a global no-go or future source-extension impossibility; or
  \item create Gate Chair, benchmark, proof, publication, release, push,
  external-action, or completed-derivation authority.
\end{itemize}

\end{document}
